\chapter{Robustness to Distribution Shift}
\label{ch:rq5}

The preceding experiments measured responses to controlled noise and local
perturbations on in-distribution ImageNet images or their token grids. Inputs
encountered outside training can differ systematically in style, content, and
layout. This chapter tests whether these structured changes reduce tokenizer
robustness by evaluating six out-of-distribution datasets.
ImageNet validation remains the in-distribution reference, and six shifted
datasets vary sampling, visual style, object presentation, imaging domain, and
document layout. The experiment measures changes in global code frequency,
positional code diversity and reconstruction quality. Its
central question is whether larger changes in code usage coincide with worse
reconstruction.

Both tokenizers were trained on natural images. Chapter~\ref{ch:rq1} evaluates
them on the ImageNet validation set as the in-distribution reference. A shifted dataset can change their token
distributions in two ways. First, it may activate codes that were absent on
ImageNet, indicating that the ImageNet baseline did not expose the tokenizer's
operational vocabulary. Second, the
active vocabulary may remain fixed while assignment probabilities change,
concentrating many image patches onto a smaller effective set of codes. Raw
active-code counts detect the first failure but not the second.

Spatial allocation provides a separate view. Because each image is represented
by a $16\times16$ grid, a code can be frequent globally yet appear in different
parts of the grid under a new domain. Similarly, the range of codes observed at
a fixed grid position may contract when a dataset has strong layout regularity.
The chapter therefore separates three questions:
\begin{enumerate}
    \item Does the global code-frequency distribution change?
    \item Does code diversity change across grid positions?
    \item Are these changes accompanied by worse reconstructions?
\end{enumerate}

\section{Experimental Setup}
\label{sec:rq5_setup}

\subsection{Dataset-Shift Protocol}
\label{subsec:rq5_dataset_protocol}

The fixed reference is the 50{,}000-image ImageNet validation set. It is
compared separately with ImageNet-V2, ImageNet-Sketch, ObjectNet, OrganAMNIST,
BloodMNIST, and RVL-CDIP, using the datasets and preprocessing described in
Chapter~\ref{ch:datasets_models_metrics}. The comparison is always ImageNet
versus one shifted dataset. Table~\ref{tab:rq5_dataset_sizes} reports the number of encoded images.
Both tokenizers produce a $16\times16$ grid, giving 256 token assignments per
image.

\begin{table}[t]
    \centering
    \small
    \begin{tabular}{llr}
        \textbf{Role} & \textbf{Dataset} & \textbf{Images} \\ \hline
        ID & ImageNet & 50{,}000 \\ \hline
        OOD & ImageNet-V2 & 10{,}000 \\
        OOD & ObjectNet & 50{,}273 \\
        OOD & ImageNet-Sketch & 50{,}889 \\
        OOD & OrganAMNIST & 57{,}187 \\
        OOD & BloodMNIST & 17{,}092 \\
        OOD & RVL-CDIP & 39{,}996 \\
    \end{tabular}
    \caption{Datasets encoded for the distribution-shift analysis. ImageNet is
    the in-distribution reference for every comparison.}
    \label{tab:rq5_dataset_sizes}
\end{table}

For each tokenizer and dataset, the experiment exports the global count of
every code and a $K\times16\times16$ tensor of code counts by grid position.
The counts estimate code frequencies across the dataset and at each grid
position.

\subsection{Global Code-Usage Shift}
\label{subsec:rq5_global_metrics}

Let $p$ be the ImageNet code-frequency distribution defined in
Section~\ref{subsec:codebook_usage_metrics}, and let $q$ be the corresponding
distribution on an OOD dataset. We define two ratios for this comparison,
active-set retention and perplexity ratio, and use Jensen--Shannon divergence
as a standard measure of distribution difference.

\textbf{Active-set retention} is defined as
\begin{equation}
    \operatorname{Ret}(p,q)=
    \frac{|\{k:p_k>0\}\cap\{k:q_k>0\}|}{|\{k:p_k>0\}|}.
    \label{eq:rq5_active_retention}
\end{equation}
It measures the fraction of codes observed on the ID dataset that also occur
on the OOD dataset. Retention equals one if all ID-active codes remain active.

Perplexity is the exponential of distribution entropy,
$P(p)=\exp(H(p))$. The \textbf{perplexity ratio} is defined as
\begin{equation}
    R_P = \frac{P(q)}{P(p)}
    \label{eq:rq5_perplexity_ratio}
\end{equation}
and measures the change in effective vocabulary size. A value near one indicates
similar spread; a value below one indicates concentration onto fewer effective
codes.

The \textbf{Jensen--Shannon divergence} measures how much the probabilities of
the same code identities differ between datasets. It is defined as:
\begin{equation}
    \operatorname{JSD}(p,q)
    = \frac{1}{2}D_{\mathrm{KL}}(p\|m)
    + \frac{1}{2}D_{\mathrm{KL}}(q\|m),
    \qquad m=\frac{p+q}{2}.
    \label{eq:rq5_jsd}
\end{equation}
Here $m$ is the midpoint distribution. JSD is symmetric and remains finite
when a code appears in only one dataset; direct KL divergence need not.
It is zero for identical distributions and at most $\log 2$ with natural
logarithms. Unlike a comparison between independently sorted rank-frequency curves, this
quantity preserves code identity: probability assigned to code $k$ on
ImageNet is compared with probability assigned to the same code $k$ on the
shifted dataset. All entropy and divergence calculations use natural
logarithms, so values are reported in nats. The comparison uses $R_P$ to
measure concentration relative to each tokenizer's own ID baseline.

These global quantities pool millions of token assignments per dataset and
are therefore less sparse than the per-code positional estimates below.
Nevertheless, the datasets have unequal image counts, so small differences,
especially for ImageNet-V2 and BloodMNIST, are not treated as formal
hypothesis-test results.

\subsection{Positional Measurements}
\label{subsec:rq5_positional_metrics}

For grid position $(u,v)$, let $p_{k\mid u,v}$ be the probability of observing
code $k$ at that position over all images in a dataset. Positional entropy is
\begin{equation}
    H_{u,v}=-\sum_{k=1}^{K}p_{k\mid u,v}\log p_{k\mid u,v}.
    \label{eq:rq5_positional_entropy}
\end{equation}
We normalize entropy by $\log K$, its theoretical maximum for a codebook of
size $K$. Both checkpoints have $K=16{,}384$. The plots report the average of
$(H^{\mathrm{shift}}_{u,v}-H^{\mathrm{ImageNet}}_{u,v})/\log K$ over all
positions and separately over border and interior positions.
Negative values mean that fewer codes, or a more concentrated mix of codes,
occur at that location under the shifted dataset.

These positional entropy estimates depend on sample size. Each position
contributes one observation per image, so an empirical entropy computed from
$N$ images cannot exceed $\log(\min(K,N))$. ImageNet-V2 has 10{,}000 images,
fewer than LlamaGen's 16{,}384 codes. Its entropy difference consequently
includes a sampling effect and cannot establish a loss of intrinsic diversity.

\subsection{Reconstruction Measurements}
\label{subsec:rq5_reconstruction_metrics}

PSNR, SSIM, LPIPS, and FID are computed between each dataset and its
reconstructions under the protocol of Chapter~\ref{ch:datasets_models_metrics}.
Paired metrics are reported consistently across rows, but their magnitudes
remain sensitive to dataset content: a homogeneous medical dataset can yield
high PSNR without being close to ImageNet. FID is used only to compare the two
tokenizers on the same dataset.
Its absolute value is not compared across datasets because each row has a
different reference distribution. FID also inherits the natural-image bias of
the Inception features on which it is based; this is particularly relevant for
the medical and document datasets.

\section{Experimental Results}
\label{sec:rq5_results}

\subsection{The Active Vocabulary Is Stable but Its Distribution Is Not}
\label{subsec:rq5_global_results}

Table~\ref{tab:rq5_global_shift} reports the global code-usage measurements.
ImageNet-V2 produces almost no global shift: JSD is 0.002 for LlamaGen and
0.0004 for VQGAN, while both perplexity ratios are approximately one.
ObjectNet also remains close to the ImageNet distribution, with JSD values of
0.020 and 0.007. ImageNet-V2 contains natural images of the same classes collected from
different sources. ObjectNet changes object presentation. Both produce small
changes in the global code-frequency distribution.

\begin{table}[t]
    \centering
    \scriptsize
    \setlength{\tabcolsep}{4pt}
    \begin{tabular}{llrrr}
        \textbf{Tokenizer} & \textbf{Shifted dataset} & \textbf{Active retention} & \textbf{JSD} & $\boldsymbol{R_P}$ \\ \hline
        LlamaGen & ImageNet-V2     & 1.000 & 0.0020 & 0.995 \\
        LlamaGen & ObjectNet       & 1.000 & 0.0200 & 0.916 \\
        LlamaGen & ImageNet-Sketch & 1.000 & 0.1531 & 0.377 \\
        LlamaGen & OrganAMNIST     & 1.000 & 0.1035 & 0.624 \\
        LlamaGen & BloodMNIST      & 0.984 & 0.2260 & 0.359 \\
        LlamaGen & RVL-CDIP        & 0.994 & 0.2510 & 0.230 \\ \hline
        VQGAN & ImageNet-V2        & 1.000 & 0.0004 & 1.001 \\
        VQGAN & ObjectNet          & 1.000 & 0.0068 & 0.962 \\
        VQGAN & ImageNet-Sketch    & 1.000 & 0.1277 & 0.452 \\
        VQGAN & OrganAMNIST        & 1.000 & 0.0424 & 0.816 \\
        VQGAN & BloodMNIST         & 0.994 & 0.1575 & 0.497 \\
        VQGAN & RVL-CDIP           & 0.999 & 0.1949 & 0.325 \\
    \end{tabular}
    \caption{Global code-usage shift relative to ImageNet. Active retention is
    the fraction of ImageNet-active codes also observed on the shifted dataset;
    $R_P$ is the shifted-to-ImageNet perplexity ratio. }
    \label{tab:rq5_global_shift}
\end{table}

The style and domain changes concentrate the token distributions. RVL-CDIP gives
the largest JSD for both tokenizers, 0.251 for LlamaGen and 0.195 for VQGAN.
BloodMNIST and ImageNet-Sketch follow. Their perplexity ratios fall to
0.359--0.377 for LlamaGen and 0.452--0.497 for VQGAN. OrganAMNIST is
intermediate. The ordering is broadly shared by the two models even though
their ImageNet baselines are very different.

Active retention alone would miss this result. LlamaGen retains at least
98.4\% of its 16{,}384 ImageNet-active codes on every dataset. VQGAN retains at
least 99.4\% of its 972 ImageNet-active codes, and exactly the same 972 codes
remain active on ImageNet-V2, ImageNet-Sketch, ObjectNet, and OrganAMNIST. The
shift therefore reallocates probability among an almost unchanged vocabulary;
it does not activate ImageNet-inactive VQGAN codes at scale or collapse
LlamaGen to a small raw set.

\begin{figure}[t]
    \centering
    \includegraphics[width=0.98\textwidth]{ch08_fig_global_shift_summary.png}
    \caption{Global response to each shifted dataset. JSD measures movement of
    probability mass over fixed code identities; $R_P$ measures concentration relative to each tokenizer's ImageNet
    baseline.}
    \label{fig:rq5_global_shift_summary}
\end{figure}

Figure~\ref{fig:rq5_global_shift_summary} shows that ImageNet-V2 and ObjectNet
remain near the ImageNet reference. The stylized, medical, and document
datasets have lower perplexity ratios, despite little
change in their raw active sets.

\subsection{Reconstruction under Distribution Shift}
\label{subsec:rq5_reconstruction_results}

Table~\ref{tab:rq5_reconstruction_shift} reports reconstruction quality on the
six shifted datasets. LlamaGen is better in every model-paired comparison: its
PSNR and SSIM are higher, while its LPIPS and FID are lower. The advantage is
not confined to the ImageNet baseline established in
Chapter~\ref{ch:rq1}.

\begin{table}[t]
    \centering
    \scriptsize
    \setlength{\tabcolsep}{4pt}
    \begin{tabular}{llrrrr}
        \textbf{Tokenizer} & \textbf{Dataset} & \textbf{PSNR} & \textbf{SSIM} & \textbf{LPIPS} & \textbf{FID} \\ \hline
        LlamaGen & ImageNet         & 20.79 & 0.558 & 0.228 & 2.19 \\
        LlamaGen & ImageNet-V2      & 20.95 & 0.572 & 0.231 & 4.34 \\
        LlamaGen & ObjectNet        & 23.88 & 0.683 & 0.208 & 2.93 \\
        LlamaGen & ImageNet-Sketch  & 20.26 & 0.752 & 0.128 & 3.69 \\
        LlamaGen & OrganAMNIST      & 26.79 & 0.805 & 0.182 & 57.16 \\
        LlamaGen & BloodMNIST       & 27.10 & 0.802 & 0.157 & 19.60 \\
        LlamaGen & RVL-CDIP         & 20.49 & 0.764 & 0.094 & 5.08 \\ \hline
        VQGAN & ImageNet            & 19.65 & 0.489 & 0.286 & 5.06 \\
        VQGAN & ImageNet-V2         & 19.83 & 0.501 & 0.289 & 7.48 \\
        VQGAN & ObjectNet           & 22.79 & 0.611 & 0.268 & 5.80 \\
        VQGAN & ImageNet-Sketch     & 18.23 & 0.696 & 0.170 & 5.98 \\
        VQGAN & OrganAMNIST         & 26.58 & 0.732 & 0.239 & 76.12 \\
        VQGAN & BloodMNIST          & 26.38 & 0.766 & 0.201 & 29.02 \\
        VQGAN & RVL-CDIP            & 18.23 & 0.699 & 0.128 & 8.02 \\
    \end{tabular}
    \caption{Reconstruction quality on ImageNet and the shifted datasets. PSNR
    and SSIM are higher-is-better; LPIPS and FID are lower-is-better. FID
    comparisons are made between tokenizers within a dataset, not across
    dataset rows.}
    \label{tab:rq5_reconstruction_shift}
\end{table}

The paired reconstruction metrics show no uniform degradation across the
OOD datasets. ImageNet-V2 is close to ImageNet for both tokenizers.
ImageNet-Sketch has lower PSNR but higher SSIM and lower LPIPS, so the
metrics do not support a single ordering of those datasets. Medical images
have higher PSNR and SSIM and lower LPIPS than ImageNet, while document
images have higher SSIM and lower LPIPS but lower PSNR. These differences
also reflect image content and complexity. The consistent comparison is
between tokenizers on each dataset: all four metrics favor LlamaGen.

Large shifts in code frequencies can therefore coexist with preserved
reconstruction quality. FID is reported for the model comparison within each
dataset; its changing reference distribution and natural-image features
limit interpretation across domains.

\subsection{Positional Diversity Contracts under Larger Shifts}
\label{subsec:rq5_entropy_results}

Figure~\ref{fig:rq5_positional_entropy} reports normalized changes in positional
entropy on a common scale. Negative values indicate a more concentrated
distribution at a typical grid position. Comparing border and interior
averages tests whether that concentration is uniform across the grid.
These averages describe spatial variation without identifying its cause.

\begin{figure}[t]
    \centering
    \includegraphics[width=\textwidth]{ch08_fig_normalized_positional_entropy.png}
    \caption{Mean positional entropy change divided by $\log K$, for all
    positions, border positions, and interior positions. The common scale
    expresses changes relative to the theoretical entropy range. Unequal
    dataset sizes still affect these estimates, especially ImageNet-V2.}
    \label{fig:rq5_positional_entropy}
\end{figure}

The ImageNet-V2 result illustrates why these maps require caution. Its global
distribution is almost unchanged, but LlamaGen's mean positional entropy is
0.68 lower. With only 10{,}000 images and 16{,}384 active codes, this decrease
is partly imposed by finite sampling. VQGAN, with only 972 active codes, shows
a much smaller mean decrease of 0.039. The larger decreases on datasets with
sample sizes close to ImageNet are less exposed to this specific mismatch, but
they remain empirical entropy estimates rather than intrinsic properties of
the encoder.

\section{Discussion}
\label{sec:rq5_discussion}

The distribution-shift experiment separates vocabulary availability from
vocabulary use. VQGAN's approximately 972 active ImageNet codes remain almost
exactly the active set on all six shifted datasets. LlamaGen likewise retains
nearly its full vocabulary. Neither tokenizer responds to a new domain by
substantially expanding or shrinking its raw active set. The measurable change
is redistribution: frequent codes become more dominant and perplexity falls.
Raw active-code count detects changes in vocabulary membership, but it misses
this redistribution within the active set.

ImageNet-V2 and ObjectNet remain close to the ImageNet token distribution.
Sketches, medical images, and documents use a more concentrated token mixture,
with RVL-CDIP producing the largest JSD. The two tokenizers show the same
qualitative JSD ranking despite their different empirical vocabulary sizes.

Changing the input domain substantially changes token frequencies for sketches,
medical images, and documents, while paired reconstruction scores show no
uniform decline. Both tokenizers can reconstruct images whose token
distribution differs from ImageNet. An autoregressive model trained on the
ID token distribution would receive a changed frequency distribution, with
potential consequences for its performance.

The global results are the main finding of this chapter. An almost unchanged
active vocabulary can hide substantial redistribution among its codes.
LlamaGen retains its reconstruction advantage across the tested datasets,
but its larger vocabulary does not prevent this redistribution.
